\section{Conclusion}

This project studies flash-loan-like interactions as an important DeFi execution pattern and focuses on three representative protocol families: Aave V3, Balancer V2, and Uniswap V2. Within this scope, we aimed to explain provider-user interaction flows, design protocol-aware identification rules, and validate those rules on real Ethereum mainnet transactions.

The main outcome of the project is a working prototype scanner rather than a purely conceptual rule list. The implemented system follows a \texttt{candidate} $\rightarrow$ \texttt{verified} $\rightarrow$ \texttt{strict} pipeline, stores results at the interaction, asset-leg, and transaction levels, and combines event-backed and trace-backed evidence to justify final labels. In the selected real-mainnet validation window, the prototype successfully detected one representative strict sample for each supported protocol, which shows that the system can produce explainable results on real transactions rather than only on synthetic examples.

Overall, the project suggests that flash-loan identification is better handled as a protocol-specific pattern-recognition problem than as a single black-box classification task. Although the current prototype is not yet a chain-wide measurement system, it already provides a clear and reproducible foundation for understanding flash-loan-like behavior in DeFi. Future work can extend this foundation with larger-scale evaluation, broader protocol coverage, and richer downstream transaction analysis.
